\documentclass{article}
\usepackage{amsmath}
\usepackage{graphicx}
\usepackage{float}

\title{Final Model Report and Business Insights}
\author{Tommy Yao}
\date{August--September 2026}


\begin{document}


\maketitle

\begin{abstract}
The following is my overall statistical reports from my final model. In addition to this, I will firstly provide
specific business/strategic insights that can be derived from the final model. My final model do have some 
undeprediction problems for high-demand observations, however, it is better than the overprediction of the
high-demand observations, and I will explain the potential reasons and opportunity cost for each of them.
\end{abstract}
\section{Business Insights: Why Overprediction of High-Demand Observations is more disadvantageous for both Driver and Companies}
\subsection{Overprediction Problems and reasons its intractability}
For Instance, if the model predicts that the demand in zone A is 650 and the demand in zone B is 250-350, but the actual demand in zone A is 500 and the actual demand in zone B is 300.
We generally Assume that location A is hotspot with demand higher than 350, and location B is more general location with demand between 5-350.
\[
D_A = 500, \qquad D_B = 300
\]
\[
\hat{D}_A = 650, \qquad \hat{D}_B = 250-350
\]
According, the company
 will reallocate drivers from zone B to zone A. 
\[
\text{Driver Reallocation:} \qquad B \rightarrow A
\]
This will result in a shortage of drivers in zone B, which will lead to longer wait times for customers in zone B and potentially lost revenue for the
company.
Thus, we have:
\[
\text{Oversupply}_A+\text{Undersupply}_B
\]
I give following reasons that why it is intractable:
\begin{enumerate}
    \item \textbf{Hurt market for two zones:}
     When the prediction is made, it has already hurt two spatial markets. 
     It is hard to do any compensation to fix this effect.


    \item \textbf{Driver become unsatisfied and picky:} 
    When Driver become idle, the driver will become more picky about orders 
    and become unstatiated because of the time consumption.


    \item  \textbf{Irreversibility in the short run:}
    The problems caused by overprediction may already be irreversible in the short run.
\end{enumerate}


\subsection{Underprediction: Problems and Correction Mechanisms}


Underprediction occurs when the actual ride demand exceeds the predicted demand. 
This may result in insufficient driver supply, longer passenger waiting times, 
higher cancellation rates, and potential revenue loss. However, compared with 
overprediction, underprediction can be partially corrected through real-time 
marketplace adjustments.

\begin{enumerate}
    \item \textbf{Real-Time Demand Detection:} 
    Uber and Lyft can observe incoming ride requests and driver availability in 
    real time. Therefore, an unexpected demand surge can be detected shortly 
    after the original forecast is generated.

    \item \textbf{Dynamic Pricing and Driver Incentives:} 
    When demand exceeds available supply, surge pricing and driver incentives 
    can encourage more drivers to become active or move toward high-demand areas, 
    partially correcting the supply shortage.

    \item \textbf{Flexible Driver Repositioning:} 
    Available drivers from nearby zones can be encouraged to reposition toward 
    the under-supplied area. This allows the platform to respond dynamically 
    without requiring the original demand forecast to be perfectly accurate.
\end{enumerate}
\end{document}